% ════════════════════════════════════════════════════════════════════════════
\section{A game-based formulation}
\label{app:game}
% ════════════════════════════════════════════════════════════════════════════

\newcommand{\CRExp}{\mathsf{CR}\text{-}\mathsf{Exp}}
\newcommand{\SzExp}{\mathsf{Sz}\text{-}\mathsf{Exp}}

The body states coercion-safety in a \emph{decision-theoretic} register: the
coercion game $\mathbf{Coerce}_{\Sigma}$ (Definition~\ref{def:game}) asks whether a
\emph{priced} adversary $\mathrm{WDY}[W,c]$ ever strictly prefers to kill. Readers
from the cryptographic tradition will expect the complementary
\emph{experiment-and-advantage} form. We give it here --- and, mirroring the two
axes the residual splits into (\S\ref{sec:residual-seizure} for funds,
\S\ref{sec:residual-dr} for the deadly race), we give \emph{two} experiments over
one shared WDY interaction: a \textbf{seizure game}, whose winning event is the
familiar ``the adversary reaches $K$,'' and a \textbf{deadly-race game}, whose
winning event is the novel ``the adversary reaches a \emph{gray} state'' --- the one
that manufactures the kill incentive (Lemma~\ref{lem:drl}), and which an adversary
can trigger even when it provably cannot learn $K$.

\paragraph{Two adversaries, and why the primed one is the worst case.}
The seizure game keeps the body's adversary $\Attacker$ (holder $\Victim$): it wants
the \emph{money}, seeks $K$, and would kill only \emph{instrumentally}, where the
priced calculus $(1-p_A)\,W>c$ makes a kill raise its take
(Definition~\ref{def:priced}). Its question is fund-safety --- \emph{does $\Attacker$
reach $K$?} The deadly-race game instead posits a \emph{hypothetical} adversary
$\Attacker'$ (holder $\Victim'$, primed to keep the two games apart) that wants the
deadly race \emph{for its own sake}: it plays to reach a gray state. This is
deliberately the worst case for the kill question. Reaching a gray state is
\emph{necessary} for any kill to pay (Lemma~\ref{lem:drl}); so if even the
DR-hunting $\Attacker'$ cannot reach one --- the coercion-resistance result below,
$\Pr[\text{win}]=0$ --- then the money-seeking $\Attacker$, who would kill only
instrumentally, never lands in a state where killing $\Victim$ pays.
\emph{$\Attacker'$ trying and failing to manufacture a deadly race is exactly what
certifies that $\Attacker$ will not run one.} Throughout, $p_A$ and $\pi$ stay
unprimed: they are the body's state functionals
(Definitions~\ref{def:racers},~\ref{def:gray}), read off the resulting position, and
each game is a different \emph{condition} on them.

\subsection{The shared interaction}

\begin{definition}[The WDY interaction]\label{def:wdy-int}
Fix a custody protocol $\Pi$, a WDY adversary $\mathcal{X}$
(Assumption~\ref{ass:power}), and the honest holder $\mathcal{Y}$.
\begin{enumerate}[nosep,leftmargin=1.7em]
  \item \textbf{Setup.} The challenger runs $\Pi$'s honest setup for $\mathcal{Y}$,
  producing the master secret $K$, the co-located device and explicit-secret state
  $D$ (vault and time-lock blobs, commitment scripts, review schedules), and the
  tacit state $T$ (the holder's recognition capability). Public parameters
  $\mathrm{pp}$, including the salt $\sigma$, are given to $\mathcal{X}$
  (clause~\ref{d:kerckhoffs}).
  \item \textbf{Interaction.} Within a single attack site, $\mathcal{X}$ queries, in
  any order:
  \begin{itemize}[nosep,leftmargin=1.5em]
    \item $\mathsf{Seize}$ --- returns the co-located device and explicit-secret
    state $D$ (clauses~\ref{d:seizure},\ref{d:explicit}).
    \item $\mathsf{Coerce}(t)$ --- for a coercion budget $t$, returns every value
    $\mathcal{Y}$ is \emph{able} to transmit under full cooperation
    (clause~\ref{d:subjugation}); the tacit state $T$ is available only as an
    \emph{evaluate-only} oracle the challenger runs on $\mathcal{Y}$'s behalf, never
    as a transmissible value (clauses~\ref{d:notacit},\ref{d:noperceptual}).
    \item $\mathsf{Compute}$ --- any computation on the feasibility line
    (clause~\ref{d:compute}): at most $2^{64}$ cheap and strictly fewer than
    $2^{64}$ memory-hard evaluations.
  \end{itemize}
  Any remote delegate is not coercible through $\mathcal{Y}$
  (clause~\ref{d:nodelegatecoerce}).
  \item \textbf{Resulting position.} Let $R$ be the resulting set of racers,
  $p_A = p_A(R)$ the probability $\mathcal{X}$ then reaches $K$
  (Definition~\ref{def:game}), and $\pi = p_A(R\setminus\{\mathcal{Y}\}) - p_A(R)$
  the holder's pivotality. Each game below is a predicate on this $(p_A,\pi)$.
\end{enumerate}
\end{definition}

\subsection{The seizure game --- fund-safety}

\begin{definition}[Seizure experiment; seizure-resistance]\label{def:szexp}
The \textbf{seizure experiment} $\SzExp_{\Pi}(\Attacker)$ runs the interaction
(Definition~\ref{def:wdy-int}) with adversary $\mathcal{X} = \Attacker$ and holder
$\mathcal{Y} = \Victim$, and outputs $1$ (``$\Attacker$ \emph{seizes}'') iff
$\Attacker$ feasibly reaches $K$ --- equivalently $p_A > 0$
(Definition~\ref{def:racers}). $\Pi$ is \textbf{seizure-resistant} if
$\Pr\!\big[\SzExp_{\Pi}(\Attacker) = 1\big] = 0$ for every WDY adversary
$\Attacker$.
\end{definition}

This is the fund-safety / seizability axis --- the game face of the \emph{floor}.

\begin{corollary}[TGPO is seizure-resistant up to the residual]\label{cor:tgpo-szsecure}
The Time-Gated Perceptual Oracle (\S\ref{sec:gw}) is seizure-resistant except on the
residual \emph{seizure} window (\S\ref{sec:residual-seizure}).
\end{corollary}
\begin{proof}
By Theorem~\ref{thm:floor}, from every seizable state save that window $K_N$ meets
the floor, so no feasible path to $K$ exists and $p_A = 0$; hence
$\SzExp$ outputs $0$. On the residual seizure window $\SzExp$ may output $1$ --- the
$o_N$-live coincidence where $\Attacker$ does reach $K$; the loss there is of
\emph{funds}, bounded as argued in \S\ref{sec:residual-seizure}.
\end{proof}

\subsection{The deadly-race game --- deadly-race-safety}

\begin{definition}[Deadly-race experiment]\label{def:crexp}
The \textbf{deadly-race experiment} $\CRExp_{\Pi}(\Attacker')$ runs the interaction
(Definition~\ref{def:wdy-int}) with adversary $\mathcal{X} = \Attacker'$ and holder
$\mathcal{Y} = \Victim'$, and outputs $1$ (``$\Attacker'$ \emph{wins}'') iff
$\pi > 0$; for the individual-custody case $R = \{\Victim'\}$ this is exactly a
\emph{gray} outcome, $0 < p_A < 1$ (Definition~\ref{def:gray}).
\end{definition}

\begin{definition}[Coercion resistance]\label{def:crsecure}
$\Pi$ is \textbf{coercion-resistant} if
$\Pr\!\big[\CRExp_{\Pi}(\Attacker') = 1\big] = 0$ for every WDY adversary
$\Attacker'$.
\end{definition}

\begin{remark}[Where the two games part --- the clean loss]
The two experiments read the same position $(p_A,\pi)$ but score it differently, and
they diverge at exactly one point. At $p_A = 0$ both output $0$ (no feasible path:
the defender clears both). On the \emph{gray} interval $0 < p_A < 1$ both output $1$
($\Attacker$ holds a path and $\Victim'$ is pivotal: the defender fails both). They
part only at \emph{full compromise} $p_A = 1$: the seizure game outputs $1$ ($K$ is
gone) while the deadly-race game outputs $0$ ($\pi = 0$ --- with the adversary
already certain to win, removing the holder changes nothing, so no kill incentive
arises). That single disagreement is the whole content of \emph{``a clean loss is
deadly-race-safe''} (Remark~\ref{rem:r2-cost}): losing the funds outright is a
seizure-game defeat but a deadly-race-game \emph{non}-event. Conversely --- the
paper's negative result restated --- an adversary that provably \emph{cannot} win the
seizure game (cannot reach $K$; e.g.\ one facing a graceful-recovery vault it cannot
open within the custody window) can still win the \emph{deadly-race} game by reaching
a gray state. The deadly-race axis is therefore \emph{not} implied by fund-safety; it
is a separate axis, exactly as \S\ref{sec:residual-seizure} and
\S\ref{sec:residual-dr} are separate. Fund-safety the construction handles on its own
terms --- \emph{on top of} removing the gray race, it minimizes the window in which
the stash is seizable at all, via the wipe-before-memory-hard step
of~\eqref{eq:orbit} and the floor (Theorem~\ref{thm:floor}).
\end{remark}

\begin{remark}[The residual window: from $0$ to $\varepsilon$]
Definition~\ref{def:crsecure} is the information-theoretic, no-gray-area form. The
residual window (\S\ref{sec:residual-dr}) replaces the exact $0$ by a bounded
$\varepsilon = \Pr[\mathsf{Corr}\cap\mathsf{Esc}]$: call $\Pi$
\textbf{$\varepsilon$-coercion-resistant} if
$\Pr[\CRExp_{\Pi}(\Attacker') = 1] \le \varepsilon$ for all $\Attacker'$. The claims
below are the $\varepsilon = 0$ statements; the residual makes them
$\varepsilon$-small --- not zero --- in the sliver argued incentive-safe there.
In plain terms, this sliver is the moment the attack coincides with a \emph{live}
final orbit point $o_N$ --- currently or imminently materialized --- and
$\Attacker'$ prevents the wipe (or reset) that would otherwise close it; only then
is transmissible material briefly within reach, which is why the window is a
bounded $\varepsilon$ rather than a standing exposure.
\end{remark}

\subsection{Consistency with the main results}

\begin{theorem}[The experiment coincides with the Criterion]\label{thm:game-nga}
$\Pi$ is coercion-resistant (Definition~\ref{def:crsecure}) if and only if
$\pi = 0$ at every reachable state --- equivalently, iff $\Pi$ satisfies the
No-Gray-Area Criterion (Criterion~\ref{prin:nga}), iff $\Pi$ is coercion-safe in
the decision-theoretic game (Definition~\ref{def:game}).
\end{theorem}
\begin{proof}
$\CRExp$ outputs $1$ exactly on a reachable state with $\pi > 0$, so
$\Pr[\text{win}] = 0$ iff no reachable state is pivotal. By
Theorem~\ref{thm:safety}, $\pi = 0$ everywhere is equivalent to coercion-safety
and to Criterion~\ref{prin:nga}. The two clearance routes of
Corollary~\ref{cor:reach} --- individual custody (\ref{r:removerace}) and a remote
un-coercible racer (\ref{r:remotewinner}) --- are exactly the two ways to force
$\Pr[\text{win}] = 0$.
\end{proof}

\begin{corollary}[TGPO is coercion-resistant]\label{cor:tgpo-crsecure}
The Time-Gated Perceptual Oracle (\S\ref{sec:gw}) is coercion-resistant.
\end{corollary}
\begin{proof}
The experiment outputs $1$ only on a \emph{gray} state ($\pi > 0$); both clear
endpoints $p_A \in \{0,1\}$ give $\pi = 0$ and hence output $0$. By
Theorem~\ref{thm:floor} every state reachable by $\mathsf{Seize}$,
$\mathsf{Coerce}$, and $\mathsf{Compute}$ leaves $K$ at the floor, so
the reachable endpoint is the \emph{clear failure} $p_A = 0$ --- the opposite
endpoint $p_A = 1$ would require feasibly reaching $K$, which the same theorem
bars. Either way $\pi = 0$, so the experiment never outputs $1$ and
$\Pr[\text{win}] = 0$. The only place $p_A$ rises above $0$ (to $1$ in the worst
instant) is the residual window (\S\ref{sec:residual-seizure}), where the tacit premise
momentarily fails; accounting for it yields $\varepsilon$-coercion resistance.
\end{proof}

\begin{remark}[Reading the classification as adversaries]
Under Definition~\ref{def:crsecure}, the classification of \S\ref{sec:classify} is
a statement about which designs admit a winning $\Attacker'$. Against a
time-lock-rescue vault advertised as \emph{individual} custody, the adversary that
calls $\mathsf{Seize}$ and then waits out the delay is winning: the rescue racer is
co-located and reachable (clause~\ref{d:seizure}), so $0 < p_A < 1$ and $\pi > 0$.
Routing rescue to a genuinely remote, un-coercible delegate removes the reachable
racer and restores $\Pr[\text{win}] = 0$ --- the honest \ref{r:remotewinner} case,
at the cost of individual custody.
\end{remark}
